% =============================================================
% integration/chapter3.tex
% Integration wrapper for Chapter 3 (FROZEN SOURCE).
%
% This file adapts the standalone Chapter3/Chapter3_Paper.tex for
% inclusion in the dissertation master document without modifying
% the frozen source.
%
% Frozen source: Chapter3/Chapter3_Paper.tex
% Standalone commands suppressed: \documentclass, \begin{document},
%   \begin{titlepage}, standalone \tableofcontents, \end{document}
% Package commands suppressed (all loaded globally in config/packages.tex).
% Local settings scoped with \begingroup...\endgroup.
%
% INTEGRATION LEDGER: see docs/03_INTEGRATION_LEDGER.md
% =============================================================

% ---- Set graphics search path for this chapter --------------
\graphicspath{%
  {Chapter3/}%
  {Chapter3/figures/}%
  {Chapter3/appendix/}%
  {Chapter3/sections/}%
}

% ---- Scope Ch3 paragraph rhythm locally ---------------------
\begingroup
\setlength{\parskip}{0.55em}
\setlength{\parindent}{0pt}

% Keep oversized standalone floats inside the dissertation page.
% Scale the completed float only when it exceeds the text height; retain
% every source row, panel, caption and note. Scoped to this essay.
\newsavebox{\DissertationFloatBox}
\makeatletter
\let\dissertation@endfloatbox\@endfloatbox
\def\@endfloatbox{%
  \dissertation@endfloatbox
  \ifdim\dimexpr\ht\@currbox+\dp\@currbox\relax>\textheight
    \typeout{DISSERTATION: fitting oversized float to page height}%
    \setbox\DissertationFloatBox\box\@currbox
    \global\setbox\@currbox\vbox{%
      \hbox to\columnwidth{\hfil
        \resizebox*{!}{0.98\textheight}{\usebox{\DissertationFloatBox}}%
      \hfil}}%
  \fi
}
\makeatother

% ---- Chapter opening (replaces standalone \begin{titlepage}) ----
\chapter[Essay 3 \textemdash{} Re-visiting the Political Economy of the Rise and Fall of the Unidad Popular: Towards a Global Political Economy Approach]%
  {Re-visiting the Political Economy of the Rise\\
   and Fall of the Unidad Popular:\\
   Towards a Global Political Economy Approach}

\vspace{0.8em}

% Abstract (verbatim from frozen source)
\begin{abstract}
This chapter evaluates competing explanations of the macroeconomic crisis of
Chile's Unidad Popular (1970--1973), examining the causal ordering between
domestic monetary expansion, inflation, the external constraint, and
central-bank accommodation. Advancing dependency theory as a research programme
in the Lakatosian sense, it develops the programme's protective belt by
specifying an empirically testable balance-sheet mechanism that links the
international currency hierarchy, access to world money, the balance-of-payments
constraint, peripheral central-bank solvency, and endogenous monetary
accommodation. Historically, the analysis distinguishes the contingent exercise
of U.S.\ hegemonic agency surrounding the breakdown of Bretton Woods, the
targeted bilateral credit restrictions confronting Chile, and the constrained
but real institutional agency of the Central Bank of Chile (BCCh). Using a
multi-frequency empirical design, the empirical strategy deploys stationary
vector-autoregressive Granger directional-precedence tests and a non-linear
Threshold Vector Autoregression where the central bank's predetermined Solvency
Growth Gap ($\Delta s_{t-1} \equiv g^F_{t-1} - g^H_{t-1}$) acts as the
transition variable. The empirical analysis reveals that in the adverse
solvency-growth state ($\Delta s_{t-1} \le -4.615\%$, reflecting conditions of
central bank insolvency), forward monetary transmission (base-money shock to
inflation) is statistically significant for 9 consecutive months with a
cumulative 12-month response of 3.55 percentage points. In contrast, reverse
accommodation (inflation shock to base-money growth) is statistically
significant for 10 consecutive months, delivering a cumulative 12-month
expansion of 6.54 percentage points. Reverse accommodation exceeds forward
transmission in both duration and cumulative magnitude, with point-wise
multiplier differences significant across 9 consecutive months
($h=2,\dots,10$). These findings indicate that monetary transmission is
state-dependent on external reserve solvency, providing evidence consistent
with an endogenous accommodation mechanism under severe balance-of-payments
constraints and challenging accounts that treat money creation as an autonomous
driver of the Chilean inflation.
\end{abstract}

\vspace{0.8em}

\noindent\textbf{Keywords:} Unidad Popular, Balance of Payments, Financial
Subordination, Solvency Ratio, Endogenous Money, Threshold Vector
Autoregression, International Currency Hierarchy, Dependency Theory.

\vspace{0.4em}
\noindent\textbf{JEL Codes:} B51, C32, E58, F33, N16.

\clearpage

% ---- Chapter body (frozen source sections) ------------------
% Use Chapter3/ as base dir so that nested \input{sections/...},
% \input{tables/...}, and \input{figures/...} paths within files
% resolve correctly relative to Chapter3/ root.
\subimport{Chapter3/}{sections/01_introduction}
\subimport{Chapter3/}{sections/02_literature_review}
\subimport{Chapter3/}{sections/03_macro_framework}
\subimport{Chapter3/}{sections/04_data_architecture}
\subimport{Chapter3/}{sections/05_historical_empirical_results}
\subimport{Chapter3/}{sections/06_discussion_conclusion}

% ---- Chapter bibliography -----------------------------------
\clearpage
{%
\setlength{\bibsep}{2.5pt}
\bibliographystyle{apalike}
\bibliography{Chapter3/references}
}

% ---- Chapter appendices -------------------------------------
\clearpage
\begin{subappendices}
\subimport{Chapter3/}{appendix/appendix_bop_levr}
\clearpage
\subimport{Chapter3/}{appendix/appendix_archival_codebook}
\clearpage
\subimport{Chapter3/}{appendix/appendix_causal_pcmci_ee1}
\clearpage
\subimport{Chapter3/}{appendix/appendix_tvar_girf_atlas}
\clearpage
\subimport{Chapter3/}{appendix/appendix_linear_var_diagnostics}
\end{subappendices}

\endgroup  % end Ch3 paragraph-rhythm scope
